% from ega1-7-fr.tex
\oldpage[I]{163}
noethérien. Alors $\mathscr{R}(X)$ est une $\mathscr{O}_X$-Algèbre
quasi-cohérente~; si en outre $X$ est réduit (ce qui est le cas lorsque $X$
est localement intègre), l'homomorphisme canonique
$\mathscr{O}_X\to\mathscr{R}(X)$ est injectif.
\end{proposition}

La question étant locale, la première assertion résulte de
(\hyperref[I.7.3.3-fr]{7.3.3})~; la seconde résulte aussitôt de
(\hyperref[I.7.2.3-fr]{7.2.3}).

\begin{env}[7.3.8]
\label{I.7.3.8-fr}
Soient $X$, $Y$ deux préschémas ayant chacun un nombre fini de composantes
irréductibles, et soit $f:X\to Y$ un morphisme dont la restriction à
l'ensemble des points génériques des composantes irréductibles de $X$ est une
surjection sur l'ensemble des points génériques des composantes irréductibles
de $Y$. Alors on a
\[
  f^*(\mathscr{R}(Y))=\mathscr{R}(X).
  \tag{7.3.8.1}\label{I.7.3.8.1-fr}
\]
\end{env}

En effet, on est ramené (en vertu de
(\hyperref[I.7.3.3-fr]{7.3.3})) au cas où $X$ et $Y$ sont irréductibles, de
points génériques $x$, $y$, avec $f(x)=y$~; donc
\[
  (f^*(\mathscr{R}(Y)))_x
  =\mathscr{O}_y\otimes_{\mathscr{O}_y}\mathscr{O}_x
  =\mathscr{O}_x
\]
(\hyperref[0.4.3.1-fr]{0, 4.3.1}), ce qui démontre (7.3.8.1) en vertu de
(\hyperref[I.7.3.5-fr]{7.3.5}).

\subsection{Faisceaux de torsion et faisceaux sans torsion.}
\label{subsection:I.7.4-fr}

\begin{env}[7.4.1]
\label{I.7.4.1-fr}
Soit $X$ un préschéma \emph{intègre}. Pour tout $\mathscr{O}_X$-Module
$\mathscr{F}$, l'homomorphisme canonique
$\mathscr{O}_X\to\mathscr{R}(X)$ définit par tensorisation un homomorphisme
(dit encore \emph{canonique})
\[
  \mathscr{F}\to
  \mathscr{F}\otimes_{\mathscr{O}_X}\mathscr{R}(X)
\]
qui, sur chaque fibre, n'est autre que l'homomorphisme $z\to z\otimes 1$ de
$\mathscr{F}_x$ dans
$\mathscr{F}_x\otimes_{\mathscr{O}_x}R(X)$. Le \emph{noyau $\mathscr{T}$}
de cet homomorphisme est un sous-$\mathscr{O}_X$-Module de $\mathscr{F}$,
appelé \emph{faisceau de torsion de $\mathscr{F}$}~; il est quasi-cohérent si
$\mathscr{F}$ est quasi-cohérent
(\hyperref[I.4.1.1-fr]{4.1.1} et \hyperref[I.7.3.6-fr]{7.3.6}). On dit que
$\mathscr{F}$ est \emph{sans torsion} si $\mathscr{T}=0$ et que
$\mathscr{F}$ est un \emph{faisceau de torsion} si
$\mathscr{T}=\mathscr{F}$. Pour tout $\mathscr{O}_X$-Module $\mathscr{F}$,
$\mathscr{F}/\mathscr{T}$ est sans torsion. On déduit de
(\hyperref[I.7.3.5-fr]{7.3.5}) que~:
\end{env}

\begin{proposition}[7.4.2]
\label{I.7.4.2-fr}
Si $X$ est un préschéma intègre, tout $\mathscr{O}_X$-Module quasi-cohérent
sans torsion $\mathscr{F}$ est isomorphe à un sous-faisceau $\mathscr{G}$
d'un faisceau simple de la forme $(\mathscr{R}(X))^{(I)}$, engendré (en tant
que $\mathscr{R}(X)$-Module) par $\mathscr{G}$.
\end{proposition}

Le cardinal de $I$ est appelé le \emph{rang de $\mathscr{F}$}~; pour tout
ouvert affine non vide $U$ de $X$, le rang de $\mathscr{F}$ est égal au rang
de $\Gamma(U,\mathscr{F})$ en tant que $\Gamma(U,\mathscr{O}_X)$-module,
comme on le voit aussitôt en considérant le point générique de $X$, contenu
dans $U$. En particulier~:

\begin{corollary}[7.4.3]
\label{I.7.4.3-fr}
Sur un préschéma intègre $X$, tout $\mathscr{O}_X$-Module quasi-cohérent sans
torsion et de rang $1$ (en particulier tout $\mathscr{O}_X$-Module
inversible) est isomorphe à un sous-$\mathscr{O}_X$-Module de
$\mathscr{R}(X)$, et réciproquement.
\end{corollary}

\begin{corollary}[7.4.4]
\label{I.7.4.4-fr}
Soient $X$ un préschéma intègre, $\mathscr{L}$, $\mathscr{L}'$ deux
$\mathscr{O}_X$-Modules sans torsion, $f$ (resp. $f'$) une section de
$\mathscr{L}$ (resp. $\mathscr{L}'$) au-dessus de $X$. Pour que
$f\otimes f'=0$, il faut et il suffit que l'une des sections $f$, $f'$ soit
nulle.
\end{corollary}

Soit $x$ le point générique de $X$~; on a par hypothèse
$(f\otimes f')_x=f_x\otimes f'_x=0$. Comme $\mathscr{L}_x$ et
$\mathscr{L}'_x$ s'identifient à des sous-$\mathscr{O}_x$-Modules du corps
$\mathscr{O}_x$, la relation précédente entraîne $f_x=0$ ou $f'_x=0$, et par
suite $f=0$ ou $f'=0$ puisque $\mathscr{L}$ et $\mathscr{L}'$ sont sans
torsion (\hyperref[I.7.3.5-fr]{7.3.5}).

\begin{proposition}[7.4.5]
\label{I.7.4.5-fr}
Soient $X$, $Y$ deux préschémas intègres, $f:X\to Y$ un morphisme dominant.
Pour tout $\mathscr{O}_X$-Module quasi-cohérent sans torsion $\mathscr{F}$,
$f_*(\mathscr{F})$ est un $\mathscr{O}_Y$-Module sans torsion.
\end{proposition}

